%% sections/03-rubric.tex

\section{The Marathon Playtest Rubric}
\label{sec:rubric}

\subsection{Design Principles}

The rubric was designed around four principles. First, \textit{evidence-basis}: every
score requires specific textual evidence from the session log; general impressions are
explicitly excluded. Second, \textit{dimension independence}: each dimension scores a
distinct quality; a session that excels in one dimension should not automatically score
higher in another. Third, \textit{forward-only amendment}: once a version is ratified,
it applies to all future sessions; prior session scores are never retroactively adjusted.
Fourth, \textit{conservative thresholds}: the binding pass threshold (56/80, 70\%)
is set below the mean of a mature session, ensuring that sessions regularly pass without
the threshold inflating scores.

\subsection{The Eight Dimensions}

\begin{table}[t]
\caption{The 8-Dimension Rubric: Question, Bands, and Key Anchor Versions}
\label{tab:dimensions}
\begin{tabular}{@{}llll@{}}
\toprule
Dim. & Core Question & 10-band Requirement & Key Amendments \\
\midrule
Eng. & Did each PC contribute? & All PCs + unscripted outcome & --- \\
MF & Did dice land honestly? & Dice compound narratively & v1.1, v1.4, v1.11 \\
Pac. & Did scenes earn their length? & Session arc visible & --- \\
CA  & Did PCs make choices that mattered? & Choices required \& shaped outcome & v1.9, v1.10 \\
MFid & Did the adventure deliver? & All symptoms + hidden payoffs & v1.2, v1.3 \\
AL  & Did emotional moments hit? & Outliving image + inter-PC chain & v1.1–v1.9, v1.12 \\
Sur & Did anything emergent happen? & 7+ surprises, 1 amendment-worthy & --- \\
TR  & Could DM run cold? & All lookups inline & --- \\
\bottomrule
\end{tabular}
\end{table}

Table~\ref{tab:dimensions} summarizes the eight dimensions. We describe the most
analytically significant dimensions in detail below.

\paragraph{Mechanical Fairness (MF).}
MF was the dimension most frequently amended (versions v1.1, v1.4, v1.11). The core
question --- did dice land honestly and outcomes feel earned? --- is necessary but
insufficient. Version v1.1 added a \textit{cross-dimensional-harmlessness anchor}: a
session mechanic that systematically damages outcomes in other dimensions (e.g., an
exhaustion mechanic that inflicts disadvantage on atmospheric-gate rolls) caps
Mechanical Fairness at 6 regardless of dice honesty. This anchor emerged from Session 1,
in which the adventure's Cold Pulse mechanic (outcome 3: DC 10 Constitution save or 1
level of exhaustion, with disadvantage on ability checks) caused two atmospheric payoffs
to fail --- not because the dice were mishandled, but because the mechanic itself was
cross-dimensionally damaging.

Version v1.4 added a \textit{cross-character mechanical compound} requirement for the
10-band: at least one instance where a PC's class or spell mechanic is applied to another
PC's roll, framed in-voice as a character action. Examples from the corpus include
Kessa (Divination Wizard) using Portent to replace Grom's Religion roll while narrating
``Vethrenn, let it be through me.'' This constitutes a mechanical compound (spell
mechanic applied cross-character) framed as in-character action (Kessa invokes her
mentor) rather than metagame resource management.

Version v1.11 extended this to \textit{Portent-as-directed-foreknowledge}: the Portent
die voiced to a named ally before the roll fires, framing the Divination Wizard's class
fantasy (the wizard sees what others cannot and gives that seeing to someone specific).
Silent substitution does not qualify; the direction to a named PC is required.

\paragraph{Atmospheric Landing (AL).}
AL received the most amendments of any dimension (v1.1, v1.2, v1.4, v1.5, v1.6, v1.7,
v1.8, v1.9, v1.12). The core insight from version v1.1 was that atmospheric moments
have both a \textit{finder} (the PC who rolls or notices the beat) and a
\textit{receiver} (the PC the beat emotionally lands on), and these may be different
PCs. Session 2 produced a canonical instance: Kessa rolled Investigation 22 and found
Mensa's unsent bread-scrap; Grom's silence after the reading was the session's emotional
impact. The per-PC two-column tracking (v1.2) formalizes this finding/receiving
distinction.

Progressive amendments to AL built a set of \textit{paths to the 10-band}: any one of
(a) a specific sensory/emotional image that will outlive the session in players' memory
with at least one inter-PC chain (v1.1), (b) a multi-scene silent-reception chain
where a PC holds an un-articulated understanding across 3+ scenes and acts on it (v1.4),
(c) an act-without-announcement where a PC acts differently as a result of a prior naming
event without declaration (v1.5), or (d) simultaneous arc-convergence where a single
naming act triggers multiple outstanding arc-completions (v1.6). By v1.12, there were
six distinct paths to the 10-band on Atmospheric Landing.

\paragraph{Character Agency (CA).}
CA received two amendments (v1.9, v1.10) that addressed a key gap: the rubric's
initial 10-band required that session outcomes depend on specific PC choices the module
anticipated but did not require. This is a valid description of individual-decision
agency, but Campaign 4 produced a different form: \textit{sustained behavioral pattern}
--- a sequence of non-actions across multiple scenes that collectively constitutes the
``choice.'' The Thorngate Watch party earned garrison trust not through any single
identifiable decision but through four sessions of consistent non-credit-seeking
behavior. No one scene was the choice; the maintenance of the pattern was.

Version v1.10 added a third path: \textit{character-register conversion}, where a PC
converts a mechanical situation into a character-authored outcome traceable to their
character sheet heuristics rather than tactical calculation. In Campaign 4 Session 3,
the PC Davan ended a combat encounter mid-round by naming the NPC's coercion aloud
(``We know about your mother'') --- a resolution that came from her honor-accounting
heuristic rather than from damage calculation.

\subsection{Amendment History}

Table~\ref{tab:amendments} summarizes all 13 rubric versions.

\begin{table*}[t]
\caption{Rubric Amendment History (v1.0–v1.12): Dimension, Change, Source Innovations}
\label{tab:amendments}
\begin{tabular}{@{}lllp{5.5cm}@{}}
\toprule
Version & Date & Dimension & Change Summary \\
\midrule
v1.0 & 2026-04-18 & --- & Initial 8-dimension rubric, 0--10 bands \\
v1.1 & 2026-04-18 & MF, AL & MF: cross-dimensional-harmlessness anchor. AL: per-PC reception + inter-PC chain \\
v1.2 & 2026-04-19 & MFid, AL & MFid: anchor-level passive-Perception fallback + recovery-path. AL: finder/receiver two-column + silent reception \\
v1.3 & 2026-04-19 & MFid & Positional-redundancy: anchor symptoms need 2 of 3 fallback paths \\
v1.4 & 2026-04-19 & MF, AL & MF: cross-character compound at 10-band. AL: multi-scene silent-reception chain at 9+ \\
v1.5 & 2026-04-20 & AL & Act-without-announcement at 9+ band: cross-session behavioral consequence \\
v1.6 & 2026-04-21 & AL & Simultaneous arc-convergence as 10-band path (three paths now) \\
v1.7 & 2026-04-21 & AL & Framework-failure arc-activation at 7+ band \\
v1.8 & 2026-04-21 & AL & Arc-activation through behavioral stance at 7+ band \\
v1.9 & 2026-04-22 & AL, CA & AL: grief-paragraph situational resonance at 7+. CA: sustained behavioral pattern at 9+ \\
v1.10 & 2026-04-22 & CA & Character-register-conversion as third path to CA 10-band \\
v1.11 & 2026-04-22 & MF & Portent-as-directed-foreknowledge as valid cross-character compound path \\
v1.12 & 2026-04-22 & AL & Intra-session restraint anchor at 9+: same-session variant of act-without-announcement \\
\bottomrule
\end{tabular}
\end{table*}

\subsection{Amendment Mechanism}

Amendments are triggered when two or more innovation log entries cluster in the same
dimension. An innovation log entry is created when the session-gate reviewer identifies
a quality outcome that existing rubric anchors cannot score --- not a failure of the
anchors, but a surplus they cannot capture. The amendment pipeline is:

\begin{enumerate}
  \item Session gate reviewer identifies a surprising outcome.
  \item Outcome is logged to the innovations log with dimension tag and scope classification.
  \item When 2+ logged entries cluster in the same dimension, an amendment proposal is drafted.
  \item Proposal is reviewed: if the described pattern is genuinely novel (not already captured
        by an existing anchor), the proposal is adopted, the rubric version is incremented, and
        constituent innovations are marked as adopted.
  \item The new anchor applies to all sessions scored after the amendment date.
\end{enumerate}

The forward-only constraint ensures that session scores are comparable within the
version-window under which they were scored. Cross-version comparison is explicit
(we report gate scores with the applicable rubric version) rather than conflated.

\subsection{Threshold Design}

The binding pass threshold is 56/80 (70\%). Sessions 1--3 are advisory only (the
rubric is not yet calibrated to the adventure format); sessions 4 onward are binding.
The 70\% threshold was set to be achievable by a well-designed session without requiring
near-perfect scores: a session scoring 7/10 on all eight dimensions (56 points) passes
exactly. The empirical mean of passing sessions in the corpus is 74.8/80, indicating
the threshold is not inflationary.
